% =====================================================================
% pertemuan-01.tex
% Pertemuan 1 - Pengenalan Kompilasi dan Arsitektur Kompilator
% Sub-CPMK 1.1, 1.2, 1.3 - Teori terapan
% =====================================================================
\documentclass[aspectratio=169,10pt,t]{beamer}
\usepackage{config/beamerthemeubb}

\UBBmeeting{1}

\begin{document}

\UBBtitleframe{1}{Pengenalan Kompilasi dan Arsitektur Kompilator}{Teori terapan}

% ---------------------------------------------------------------------
\begin{frame}{Capaian Pertemuan}
  \textbf{CPMK induk:} CPMK 01 --- Menjelaskan tahapan kompilasi
  (\textit{lexical}, \textit{syntax}, \textit{semantic analysis}).

  \vspace{3pt}
  \begin{block}{Sub-CPMK yang dinilai pada pertemuan ini}
    \footnotesize
    \begin{itemize}
      \setlength{\itemsep}{1pt}
      \item \textbf{Sub-CPMK 1.1:} Menjelaskan perbedaan interpreter dan
            compiler serta peran assembler, interpreter, dan compiler
            sebagai translator, dengan menyebut contoh masukan dan
            keluaran masing-masing.
      \item \textbf{Sub-CPMK 1.2:} Menjelaskan fase-fase kompilator
            (analisis leksikal, sintaksis, semantik, dan pembangkitan
            kode antara) beserta peran tabel simbol dan penanganan
            kesalahan.
      \item \textbf{Sub-CPMK 1.3:} Membuat tabel penelusuran fase
            kompilasi untuk satu program C sederhana, yang memuat
            keluaran antara berupa token, pohon sintaks, dan kode antara.
    \end{itemize}
  \end{block}
\end{frame}

% ---------------------------------------------------------------------
\begin{frame}{Agenda Pertemuan (150 menit)}
  \small
  \begin{tabular}{@{}p{1.4cm}p{12.3cm}@{}}
    \toprule
    \textbf{Menit} & \textbf{Kegiatan} \\
    \midrule
    10 & Pembuka, orientasi capaian, dan peta proyek mini-compiler \\
    30 & Keluarga penerjemah; kompilator dibandingkan interpreter \\
    35 & Fase analisis dan sintesis; tabel simbol dan penanganan kesalahan \\
    30 & Penelusuran satu program C melalui token, pohon, dan kode antara \\
    30 & Latihan terstruktur: mengisi tabel penelusuran fase \\
    15 & Diskusi hasil latihan, refleksi, dan penugasan \\
    \bottomrule
  \end{tabular}

  \vspace{6pt}
  {\scriptsize\color{ubbblue!75}Pertemuan teori terapan: ceramah singkat
  diikuti latihan terstruktur dan pemeriksaan bersama.}
\end{frame}

% ---------------------------------------------------------------------
\begin{frame}{Posisi Pertemuan pada Rantai Proyek}
  \UBBchain{0}

  \vspace{2pt}
  \begin{itemize}
    \item Proyek semester: \textit{front-end} mini-compiler untuk subset
          bahasa C, dibangun bertahap.
    \item Pertemuan ini memberi peta seluruh rantai: lexer, grammar,
          parser, semantik, \textit{symbol table}, dan AST.
    \item Setiap pertemuan berikutnya menambah satu simpul rantai itu.
  \end{itemize}

  \vspace{2pt}
  \begin{exampleblock}{Milestone M0}
    Menyiapkan repositori, spesifikasi Subset-C v0.1, dan program contoh
    sebagai fondasi Praktikum 1.
  \end{exampleblock}
\end{frame}

% ---------------------------------------------------------------------
\begin{frame}{Keluarga Penerjemah}
  \begin{block}{Penerjemah (\eng{translator})}
    Program yang mengubah program sumber ke bahasa lain; keluarannya
    dapat berupa assembly, kode mesin, \textit{bytecode}, atau bahasa
    tingkat tinggi lain.
  \end{block}

  \vspace{4pt}
  \small
  \begin{tabular}{@{}p{3.1cm}p{4.6cm}p{5.7cm}@{}}
    \toprule
    \textbf{Jenis} & \textbf{Masukan} & \textbf{Keluaran} \\
    \midrule
    Assembler & Bahasa assembly & Kode mesin \\
    Kompilator & Bahasa tingkat tinggi & Assembly atau kode mesin \\
    Interpreter & Program sumber & Aksi program, tanpa kode target tetap \\
    \textit{Cross-compiler} & Program untuk platform lain & Kode untuk platform target \\
    \bottomrule
  \end{tabular}

  \vspace{4pt}
  {\footnotesize Kompilator dan interpreter sama-sama membaca program
  sumber, tetapi model eksekusinya berbeda.}
\end{frame}

% ---------------------------------------------------------------------
\begin{frame}{Kompilator dibandingkan Interpreter}
  \small
  \begin{tabular}{@{}p{3.1cm}p{5.2cm}p{5.2cm}@{}}
    \toprule
    \textbf{Aspek} & \textbf{Kompilator} & \textbf{Interpreter} \\
    \midrule
    Keluaran & Kode target yang berdiri sendiri & Tidak ada kode target permanen \\
    Waktu eksekusi & Cepat setelah kompilasi & Lebih lambat, diterjemahkan saat jalan \\
    Waktu mulai & Lebih lama karena kompilasi & Lebih cepat \\
    Portabilitas & Kode target terikat platform & Selama interpreter tersedia \\
    Analisis & Menyeluruh sebelum eksekusi & Terbatas pada saat eksekusi \\
    Contoh & GCC, Clang, javac & Python, Ruby, JavaScript \\
    \bottomrule
  \end{tabular}

  \vspace{5pt}
  \begin{exampleblock}{Batas yang tidak selalu tegas}
    Java mengompilasi sumber ke \textit{bytecode}, lalu mesin virtual
    menjalankannya dan dapat melakukan kompilasi \textit{just-in-time}.
  \end{exampleblock}
\end{frame}

% ---------------------------------------------------------------------
\begin{frame}{Cross-Compiler dan Bootstrapping}
  \begin{columns}[T,onlytextwidth]
    \begin{column}{0.48\textwidth}
      \begin{block}{\textit{Cross-compiler}}
        \footnotesize
        Berjalan di satu platform, menghasilkan kode untuk platform lain.
        Contoh: kompilasi untuk mikrokontroler dari komputer desktop.
        \begin{itemize}
          \item wajib tahu konvensi pemanggilan;
          \item ukuran kata dan tata letak memori target;
          \item ketersediaan pustaka target.
        \end{itemize}
      \end{block}
    \end{column}
    \begin{column}{0.48\textwidth}
      \begin{block}{Bootstrapping}
        \footnotesize
        Membangun kompilator suatu bahasa dengan kompilator bahasa itu
        sendiri.
        \begin{itemize}
          \item dimulai dari kompilator sederhana dalam bahasa lain;
          \item ditingkatkan bertahap sampai mandiri;
          \item GCC dan Clang memiliki sejarah bootstrapping panjang.
        \end{itemize}
      \end{block}
    \end{column}
  \end{columns}

  \vspace{4pt}
  {\footnotesize\color{ubbblue!80}Istilah \eng{compiler} berasal dari
  gagasan mengompilasi kumpulan rutin menjadi satu program; kini
  kompilator juga menganalisis, mengoptimasi, dan membangkitkan kode.}
\end{frame}

% ---------------------------------------------------------------------
\begin{frame}{Fase Analisis dan Sintesis}
  \begin{center}
  \begin{tikzpicture}[
      ph/.style={draw=ubbblue!40,fill=ubbblue!7,rounded corners=3pt,
        minimum width=2.30cm,minimum height=0.82cm,align=center,
        font=\scriptsize,text=ubbblue,inner sep=2pt},
      data/.style={ph,draw=ubbgold,fill=ubbgold!22,line width=1pt},
      ar/.style={-{Stealth[length=1.8mm]},draw=ubbblue!50,line width=0.8pt}
    ]
    \node[ph] (src) {Program sumber};
    \node[ph,right=0.22cm of src] (lex) {Analisis leksikal};
    \node[ph,right=0.22cm of lex] (syn) {Analisis sintaksis};
    \node[ph,right=0.22cm of syn] (sem) {Analisis semantik};
    \node[data,right=0.22cm of sem] (ir) {IR};
    \node[ph,below=0.62cm of ir] (opt) {Optimasi};
    \node[ph,left=0.22cm of opt] (gen) {Pembangkitan kode};
    \node[ph,left=0.22cm of gen] (tgt) {Program target};
    \draw[ar] (src) -- (lex);
    \draw[ar] (lex) -- (syn);
    \draw[ar] (syn) -- (sem);
    \draw[ar] (sem) -- (ir);
    \draw[ar] (ir) -- (opt);
    \draw[ar] (opt) -- (gen);
    \draw[ar] (gen) -- (tgt);
    \node[font=\scriptsize\bfseries,text=ubbteal,below=0.10cm of syn]
      {Analisis (\eng{front-end})};
    \node[font=\scriptsize\bfseries,text=ubbteal,below=0.10cm of gen]
      {Sintesis (\eng{back-end})};
  \end{tikzpicture}
  \end{center}

  \vspace{2pt}
  \begin{itemize}
    \item Analisis memahami struktur program sumber.
    \item Sintesis menghasilkan program target.
    \item \textit{Intermediate representation} (IR) menjadi batas keduanya.
  \end{itemize}
\end{frame}

% ---------------------------------------------------------------------
\begin{frame}{Analisis Leksikal}
  \begin{itemize}
    \item Memindai karakter demi karakter dari program sumber.
    \item Mengelompokkan karakter menjadi token: jenis dan atribut.
    \item Membuang spasi, baris baru, dan komentar, tetapi menyimpan
          lokasi untuk diagnostik.
    \item Keluaran: aliran token yang siap diproses parser.
  \end{itemize}

  \vspace{4pt}
  \begin{exampleblock}{Contoh}
    \texttt{int a = 5;} menghasilkan aliran
    \texttt{INT}, \texttt{ID(a)}, \texttt{ASSIGN},
    \texttt{INT\_LIT(5)}, \texttt{SEMI}.
  \end{exampleblock}

  \vspace{2pt}
  {\footnotesize Token \texttt{INT\_LIT} membawa atribut nilai 5;
  token \texttt{ID} membawa nama \texttt{a}.}
\end{frame}

% ---------------------------------------------------------------------
\begin{frame}{Analisis Sintaksis}
  \begin{columns}[T,onlytextwidth]
    \begin{column}{0.52\textwidth}
      \begin{itemize}
        \item Memeriksa apakah aliran token sesuai tata bahasa.
        \item Menyusun token menjadi pohon urai atau AST.
        \item Parser dapat mengisi tabel simbol awal.
        \item Kesalahan sintaksis dilaporkan dan dipulihkan.
      \end{itemize}
      \vspace{3pt}
      {\footnotesize Keluaran fase ini adalah pohon sintaks yang
      hubungan induk--anaknya sesuai aturan tata bahasa.}
    \end{column}
    \begin{column}{0.44\textwidth}
      \centering
      \begin{tikzpicture}[
          nd/.style={draw=ubbblue!40,fill=ubbblue!7,rounded corners=2pt,
            minimum width=0.9cm,minimum height=0.5cm,font=\scriptsize},
          ar/.style={-{Stealth[length=1.4mm]},draw=ubbblue!50}
        ]
        \node[nd] (plus) at (0,0) {\texttt{+}};
        \node[nd] (a) at (-1.3,-0.9) {\texttt{a}};
        \node[nd] (mul) at (1.1,-0.9) {\texttt{*}};
        \node[nd] (b) at (0.2,-1.8) {\texttt{b}};
        \node[nd] (two) at (2.0,-1.8) {\texttt{2}};
        \draw[ar] (plus) -- (a);
        \draw[ar] (plus) -- (mul);
        \draw[ar] (mul) -- (b);
        \draw[ar] (mul) -- (two);
      \end{tikzpicture}
      \par\vspace{2pt}
      {\scriptsize Pohon untuk \texttt{a + b * 2}}
    \end{column}
  \end{columns}
\end{frame}

% ---------------------------------------------------------------------
\begin{frame}{Analisis Semantik}
  \begin{itemize}
    \item Memeriksa makna program, bukan sekadar bentuknya.
    \item Pertanyaan yang dijawab: apakah identifier sudah
          dideklarasikan, apakah tipe operand cocok, apakah fungsi
          dipanggil dengan jumlah argumen yang benar.
    \item Menambahkan anotasi tipe dan referensi simbol pada AST.
    \item Keluaran: AST teranotasi yang siap diterjemahkan ke kode antara.
  \end{itemize}

  \vspace{4pt}
  \begin{exampleblock}{Contoh anotasi}
    \texttt{a : int}, \texttt{b : int}, \texttt{a + b : int};
    hasil \texttt{print} bertipe \texttt{void}; \texttt{return 0}
    bertipe \texttt{int}.
  \end{exampleblock}
\end{frame}

% ---------------------------------------------------------------------
\begin{frame}{Sintesis: Optimasi dan Pembangkitan Kode}
  \begin{columns}[T,onlytextwidth]
    \begin{column}{0.48\textwidth}
      \begin{block}{Optimasi}
        \footnotesize
        \begin{itemize}
          \item lokal: dalam satu blok dasar;
          \item global: lintas blok;
          \item masukan dan keluaran berupa kode antara yang lebih
                efisien.
        \end{itemize}
      \end{block}
    \end{column}
    \begin{column}{0.48\textwidth}
      \begin{block}{Pembangkitan kode}
        \footnotesize
        \begin{itemize}
          \item memilih instruksi target;
          \item mengalokasikan register;
          \item menghasilkan assembly atau kode mesin.
        \end{itemize}
      \end{block}
    \end{column}
  \end{columns}

  \vspace{5pt}
  {\footnotesize\color{ubbblue!80}IR memisahkan pekerjaan memahami
  program (analisis) dari pekerjaan menghasilkan program target
  (sintesis), sehingga kedua bagian dapat dikembangkan terpisah.}
\end{frame}

% ---------------------------------------------------------------------
\begin{frame}{Masukan dan Keluaran Setiap Fase}
  \small
  \begin{tabular}{@{}p{3.1cm}p{5.2cm}p{5.2cm}@{}}
    \toprule
    \textbf{Fase} & \textbf{Masukan} & \textbf{Keluaran} \\
    \midrule
    Leksikal & Karakter sumber & Aliran token \\
    Sintaksis & Aliran token & Pohon urai atau AST \\
    Semantik & AST dan tabel simbol & AST teranotasi \\
    IR & AST teranotasi & Kode antara \\
    Optimasi & Kode antara & Kode antara yang lebih efisien \\
    Pembangkitan kode & Kode antara & Assembly atau kode mesin \\
    \bottomrule
  \end{tabular}

  \vspace{4pt}
  {\footnotesize Setiap fase mempertahankan atau memperkaya informasi
  dari fase sebelumnya.}
\end{frame}

% ---------------------------------------------------------------------
\begin{frame}{Tabel Simbol}
  \begin{block}{Tabel simbol (\eng{symbol table})}
    Struktur data yang memetakan identifier ke tipe, scope, alamat
    memori, dan kategori (variabel, fungsi, parameter).
  \end{block}

  \begin{itemize}
    \item Operasi utama: \texttt{insert}, \texttt{lookup}, \texttt{delete}.
    \item Scope bersarang diimplementasikan sebagai stack tabel atau
          tabel dengan rantai scope.
    \item Informasi yang disimpan: nama, tipe, ukuran, alamat, dan
          rujukan ke deklarasi.
  \end{itemize}

  \vspace{2pt}
  {\footnotesize\color{ubbblue!80}Tabel simbol dipakai parser, analisis
  semantik, dan pembangkitan kode; karena itu disebut pusat informasi
  kompilator.}
\end{frame}

% ---------------------------------------------------------------------
\begin{frame}{Jenis Kesalahan dan Pesan yang Informatif}
  \small
  \begin{tabular}{@{}p{2.6cm}p{4.4cm}p{6.6cm}@{}}
    \toprule
    \textbf{Jenis} & \textbf{Contoh} & \textbf{Pesan informatif} \\
    \midrule
    Leksikal & karakter \texttt{@} tidak dikenal & baris 4, kolom 12: karakter \texttt{@} tidak dikenal \\
    Sintaksis & titik koma hilang & baris 7: diharapkan \texttt{;} setelah ekspresi \\
    Semantik & tipe tidak cocok & baris 9: tidak dapat menetapkan \texttt{float} ke \texttt{int} tanpa \textit{cast} \\
    \textit{Run-time} & pembagian nol & baris 11: pembagian dengan nol saat eksekusi \\
    \bottomrule
  \end{tabular}

  \vspace{5pt}
  \begin{block}{Pemulihan kesalahan}
    \footnotesize
    Kompilator yang baik tidak berhenti pada kesalahan pertama. Strategi
    umum: \textit{panic mode}, \textit{phrase-level recovery}, dan
    \textit{error production}. Pesan memuat lokasi, penyebab, dan saran
    perbaikan.
  \end{block}
\end{frame}

% ---------------------------------------------------------------------
\begin{frame}[fragile]{Contoh: Scope Bersarang pada Tabel Simbol}
  \begin{columns}[T,onlytextwidth]
    \begin{column}{0.50\textwidth}
      \begin{lstlisting}[language=C,numbers=none,caption={Dua deklarasi \texttt{x} pada scope berbeda},label={lst:p1-shadow}]
int x = 1;
{
    int x = 2;
    print(x);
}
print(x);
      \end{lstlisting}
    \end{column}
    \begin{column}{0.46\textwidth}
      \footnotesize
      \begin{itemize}
        \setlength{\itemsep}{1pt}
        \item[1.] Masuk blok: tabel scope dalam didorong ke stack.
        \item[2.] \texttt{int x = 2} menutupi (\eng{shadow}) \texttt{x} luar.
        \item[3.] \texttt{print(x)} di dalam blok memakai \texttt{x = 2}.
        \item[4.] Blok selesai: tabel scope dalam dibuang.
        \item[5.] \texttt{print(x)} terakhir memakai \texttt{x = 1}.
      \end{itemize}
      \vspace{2pt}
      {\scriptsize\color{ubbblue!75}Perilaku ini diimplementasikan pada
      Pertemuan 14 (Bab XV).}
    \end{column}
  \end{columns}
\end{frame}

% ---------------------------------------------------------------------
\begin{frame}[fragile]{Program Uji dan Aliran Tokennya}
  \begin{columns}[T,onlytextwidth]
    \begin{column}{0.44\textwidth}
      \lstinputlisting[language=C,firstline=1,lastline=9,
        caption={Program \file{hello.sc} dari proyek Subset-C},
        label={lst:p1-hello}]{../buku_ajar/code/subsetc/examples/hello.sc}
    \end{column}
    \begin{column}{0.53\textwidth}
      \footnotesize
      \textbf{Aliran token}
      \par\vspace{2pt}
      {\ttfamily\scriptsize
      INT\ \textperiodcentered\ ID(main)\ \textperiodcentered\ LPAREN\
      \textperiodcentered\ RPAREN\ \textperiodcentered\ LBRACE\
      \textperiodcentered\ INT\ \textperiodcentered\ ID(a)\
      \textperiodcentered\ ASSIGN\ \textperiodcentered\ INT\_LIT(5)\
      \textperiodcentered\ SEMI\ \textperiodcentered\ INT\
      \textperiodcentered\ ID(b)\ \textperiodcentered\ ASSIGN\
      \textperiodcentered\ INT\_LIT(7)\ \textperiodcentered\ SEMI\
      \textperiodcentered\ INT\ \textperiodcentered\ ID(hasil)\
      \textperiodcentered\ ASSIGN\ \textperiodcentered\ ID(a)\
      \textperiodcentered\ PLUS\ \textperiodcentered\ ID(b)\
      \textperiodcentered\ SEMI\ \textperiodcentered\ PRINT\
      \textperiodcentered\ LPAREN\ \textperiodcentered\ ID(hasil)\
      \textperiodcentered\ RPAREN\ \textperiodcentered\ SEMI\
      \textperiodcentered\ RETURN\ \textperiodcentered\ INT\_LIT(0)\
      \textperiodcentered\ SEMI\ \textperiodcentered\ RBRACE\
      \textperiodcentered\ EOF}
      \par\vspace{4pt}
      {\scriptsize\color{ubbblue!75}Setiap token membawa jenis, lexeme,
      dan lokasi (baris, kolom).}
    \end{column}
  \end{columns}
\end{frame}

% ---------------------------------------------------------------------
\begin{frame}{Tabel Penelusuran Fase untuk \file{hello.sc}}
  \scriptsize
  \begin{tabular}{@{}p{2.05cm}p{12.10cm}@{}}
    \toprule
    \textbf{Fase} & \textbf{Keluaran untuk program contoh} \\
    \midrule
    Leksikal & 34 token berurutan: \texttt{INT ID(main) LPAREN RPAREN
      LBRACE INT ID(a) ASSIGN INT\_LIT(5) SEMI} \ldots{}
      \texttt{RBRACE EOF} \\
    Sintaksis & \texttt{Program} $\rightarrow$ \texttt{DeklarasiFungsi(main)}
      $\rightarrow$ \texttt{Blok}: deklarasi \texttt{a}, \texttt{b},
      \texttt{hasil}; pernyataan \texttt{print(hasil)}; pernyataan
      \texttt{return 0} \\
    Semantik & \texttt{a:int}, \texttt{b:int}, \texttt{hasil:int};
      \texttt{a + b:int}; \texttt{print} menerima satu argumen
      bertipe \texttt{int}; hasil \texttt{print} bertipe \texttt{void} \\
    Kode antara & \texttt{t1 = a + b}; \texttt{hasil = t1};
      \texttt{param hasil}; \texttt{call print, 1};
      \texttt{t2 = 0}; \texttt{return t2} \\
    \bottomrule
  \end{tabular}

  \vspace{5pt}
  {\scriptsize\color{ubbblue!75}Setiap sel merujuk baris program sumber,
  bukan istilah umum. Tabel seperti ini menjadi artefak Sub-CPMK 1.3.}
\end{frame}

% ---------------------------------------------------------------------
\begin{frame}{Contoh Berpenyelesaian: Pohon ke Kode Antara}
  \begin{columns}[T,onlytextwidth]
    \begin{column}{0.46\textwidth}
      \centering
      \begin{tikzpicture}[
          nd/.style={draw=ubbblue!40,fill=ubbblue!7,rounded corners=2pt,
            minimum width=1.0cm,minimum height=0.55cm,font=\scriptsize},
          ar/.style={-{Stealth[length=1.4mm]},draw=ubbblue!50}
        ]
        \node[nd] (plus) at (0,0) {\texttt{+}};
        \node[nd] (a) at (-1.5,-1.0) {\texttt{a}};
        \node[nd] (mul) at (1.1,-1.0) {\texttt{*}};
        \node[nd] (b) at (0.1,-2.0) {\texttt{b}};
        \node[nd] (two) at (2.1,-2.0) {\texttt{2}};
        \draw[ar] (plus) -- (a);
        \draw[ar] (plus) -- (mul);
        \draw[ar] (mul) -- (b);
        \draw[ar] (mul) -- (two);
      \end{tikzpicture}
      \par {\scriptsize Pohon sintaks \texttt{a + b * 2}}
    \end{column}
    \begin{column}{0.50\textwidth}
      \begin{block}{Kode antara tiga alamat}
        \centering
        \texttt{t1 = b * 2}\\
        \texttt{t2 = a + t1}
      \end{block}
      \begin{itemize}
        \item Perkalian lebih dalam pada pohon, jadi dikerjakan lebih dulu.
        \item \texttt{t1} dan \texttt{t2} adalah variabel sementara.
        \item Urutan instruksi mencerminkan urutan evaluasi pohon.
      \end{itemize}
    \end{column}
  \end{columns}
\end{frame}

% ---------------------------------------------------------------------
\begin{frame}{Kegiatan Mahasiswa di Kelas}
  \begin{block}{Latihan terstruktur: tabel penelusuran fase (Sub-CPMK 1.3)}
    \footnotesize
    \begin{enumerate}
      \setlength{\itemsep}{1pt}
      \item Pilih satu program C maksimal sepuluh baris, misalnya dari
            \file{examples/}.
      \item Tuliskan aliran token beserta baris dan kolomnya.
      \item Gambarkan pohon sintaks untuk program tersebut.
      \item Tuliskan kode antara tiga alamat.
      \item Isi tabel dengan kolom fase, masukan, keluaran, dan bukti.
      \item Tukar tabel dengan teman dan periksa satu sama lain.
    \end{enumerate}
  \end{block}

  \vspace{4pt}
  \begin{columns}[T,onlytextwidth]
    \begin{column}{0.48\textwidth}
      {\footnotesize\textbf{Berkas uji:} \file{examples/hello.sc},
      \file{examples/arith.sc}.}
    \end{column}
    \begin{column}{0.48\textwidth}
      {\footnotesize\textbf{Luaran:} tabel penelusuran fase yang dapat
      diperiksa baris per baris.}
    \end{column}
  \end{columns}
\end{frame}

% ---------------------------------------------------------------------
\begin{frame}{Tugas dan Artefak}
  \begin{block}{Artefak Sub-CPMK 1.3}
    Tabel penelusuran fase untuk satu program C sederhana, memuat
    keluaran antara berupa \textbf{token}, \textbf{pohon sintaks}, dan
    \textbf{kode antara}.
  \end{block}

  \small
  \begin{itemize}
    \item Program maksimal sepuluh baris; setiap sel merujuk baris
          program sumber.
    \item Pohon sintaks konsisten dengan aliran token.
    \item Kode antara konsisten dengan pohon sintaks.
  \end{itemize}

  \vspace{3pt}
  \begin{block}{Tambahan proyek M0}
    \footnotesize
    Repositori dengan struktur \file{spec}, \file{include}, \file{src},
    \file{examples}, \file{tests}; spesifikasi v0.1 terbaca; minimal tiga
    program contoh; satu halaman catatan keputusan desain.
  \end{block}

  \vspace{2pt}
  {\footnotesize\color{ubbblue!75}M0 menjadi prasyarat Praktikum 1;
  tabel penelusuran menjadi bukti Sub-CPMK 1.3 pada UTS Praktik.}
\end{frame}

% ---------------------------------------------------------------------
\begin{frame}{Rangkuman dan Jembatan}
  \begin{itemize}
    \item Penerjemah mencakup assembler, kompilator, interpreter, dan
          \textit{cross-compiler}; masing-masing punya masukan dan
          keluaran berbeda.
    \item Kompilator bekerja dalam dua bagian: analisis (front-end) dan
          sintesis (back-end), dipisahkan oleh IR.
    \item Fase leksikal menghasilkan token, sintaksis menghasilkan pohon,
          semantik menambahkan tipe, dan IR menghasilkan kode antara.
    \item Tabel simbol adalah pusat informasi; pesan kesalahan yang baik
          menyebut lokasi dan penyebab.
    \item Tabel penelusuran fase membuktikan bahwa setiap keluaran dapat
          dilacak ke program sumber.
  \end{itemize}

  \vspace{4pt}
  \begin{exampleblock}{Pertemuan berikutnya}
    Pertemuan 2 memperdalam landasan formal: teori bahasa, grammar, dan
    hierarki Chomsky. Landasan ini menjelaskan mengapa lexer cukup
    memakai otomata hingga, sedangkan parser memerlukan otomata pushdown.
  \end{exampleblock}
\end{frame}

\end{document}
